\section{Appendix}
%\clearpage
%\onecolumn

\begin{figure*}[h]
	\begin{center}
		\begin{minipage}[t]{0.8\textwidth}
			%\hfill
			\subfigure{{\includegraphics[width=\columnwidth]{fig/expressive-0.eps}}}
			\hfill
		\end{minipage}
	\end{center}
	\caption{Computation of embeddings %of $u$ and $v$
		and \dual} 
	\label{fig:anti-expressive-0}
	%\vspace{-4ex}
\end{figure*}

\eat{
\begin{figure}
	%	\vspace{-6.7ex}
    \centerline{\includegraphics[width=0.8\columnwidth]{fig/query-example-2.eps}}
	\centering
	\vspace{-1.4ex}
	\caption{Example of complex query}
	\label{fig:query}
	\vspace{-1.4ex}
\end{figure}
}

\eat{
\begin{figure}
	%	\vspace{-6.7ex}
	\centerline{\includegraphics[width=0.8\columnwidth]{fig/anti-symmetry.eps}}
	\centering
	\vspace{-1.4ex}
	\caption{Counter example of anti-symmetry} 
	\label{fig:anti-anti-symmetry}
	\vspace{-4ex}
\end{figure}
}

\eat{
\subsection{Analysis of the loss function}
We show that when $M_1$ is sufficient large than 
$M_0$, \HGIN is more accurate.


\subsection{Gaps between embeddings of $Q$ and $G$}

\begin{equation*}\begin{aligned}
		P_{L} & \triangleq\quad\|T^Q_{L}(W_{1},\mathds{W})-T^G_{L}(W_{1},\mathds{W})\|_c \\
		& =\quad\bigg\|\phi(W_{1}x^Q_{L}+\sum_{r\in R} W_r \underbrace{\rho(\sum_{j\in 
				C^r(x^Q_{L})}g(T_{L-1,j}(W_{1},\mathds{W}))}_{\overset{\Delta}
			{\operatorname*{\operatorname*{=}}}R(W_{1},\mathds{W},x^Q_{L},r)}) \\
		& -\quad\phi(W_1x^G_L+\sum_{r\in R} W_r\rho(\sum_{j\in C^r(x^G_L)}g(T_{L-1,j}
		(W_1,
		\mathds{W}))))\bigg\|_c \\
		& =\quad C_{\phi}\bigg\|W_{1}(x^Q_{L}-x^G_{L}) + \sum_{r\in R} W_r(R(W_{1},
		\mathds{W},x^Q_{L},r)-R(W_{1},\mathds{W},x^G_{L},r))\bigg\|_c\\
\end{aligned}\end{equation*}


\stab
(1) When there exists a match of $Q$ in $G$, 
we need to consider the following two cases:
either both $Q$ and $G$ contain the center vertex
at level $L$.

\stab
(a) If $Q$ and $G$ contain the center vertices.

\begin{equation*}\begin{aligned}
		\|&R(W_1,\mathds{W},x^Q_L,r)-R(W_1,\mathds{W},x^G_L,r)\|_2 \\
		&\leq \quad C_\rho\bigg\| - \sum_{j\in C^r(x_L)\land j \neq x_L}
		W^j_R\bigg(g_0(T^G_{L-1,j}(W_1,\mathds{W})) \bigg) \\
		&+\sum_{j\in C^r(x_L)\land j \neq x_L}W^j_R\bigg(g_0(T^Q_{L-1,j}(W_1,
		\mathds{W})) -g_0(T^G_{L-1,j}(W_1,\mathds{W}))\bigg) \\
		&+ \bigg(\mathds{1}(x^Q_L\in C(x^Q_L))g_1(T^Q_{L-1,j}(W_1,\mathds{W})) 
		-
		\mathds{1}(x^G_L\in 
		C(x^G_L))g_1(T^G_{L-1,j}(W_1,\mathds{W}))\bigg)\bigg\|_c \\
		&=\quad C_\rho\bigg\|-\sum_{j\in C^r(x_L)\land j \neq x_L}
		W^j_R\bigg(g_0(T^Q_{L-1,j}(W_1,W_c)) \bigg) \\
		&+\sum_{j\in C^r(x_L)\land j \neq x_L}W^j_R\bigg(g_0(T^Q_{L-1,j}(W_1,
		\mathds{W})) -g_0(T^G_{L-1,j}(W_1,\mathds{W}))\bigg) \\
		&+ g_1(T^Q_{L-1,j}(W_1,\mathds{W})) -g_1(T^G_{L-1,j}(W_1,\mathds{W}))
		\bigg\|_c 
		\\
		&=\quad C_\rho\bigg\|g_1(T^Q_{L-1,j}(W_1,\mathds{W})) -g_1(T^G_{L-1,j}
		(W_1,
		\mathds{W}))-\sum_{j\in C^r(x_L)\land j \neq 
			x_L}W^j_R\bigg(g_0(T^Q_{L-1,j}(W_1,W_c)) 
		\bigg) \\
		&+\sum_{j\in C^r(x_L)\land j \neq x_L}W^j_R\bigg(g_0(T^Q_{L-1,j}(W_1,
		\mathds{W})) -
		g_0(T^G_{L-1,j}(W_1,\mathds{W}))\bigg)\bigg\|_c \\
		&\leq \quad C_\rho\bigg\|g_1(T^Q_{L-1,j}(W_1,\mathds{W})) -
		g_1(T^G_{L-1,j}(W_1,
		\mathds{W}))-\sum_{j\in C^r(x_L)\land j \neq 
			x_L}W^j_R\bigg(g_0(T^Q_{L-1,j}(W_1,W_c)) 
		\bigg)\bigg\|_c \\
		&+C_\rho\bigg\|\sum_{j\in C^r(x_L)\land j \neq x_L}W^j_R\bigg(g_0(T^Q_{L-1,j}
		(W_1,
		\mathds{W})) -g_0(T^G_{L-1,j}(W_1,\mathds{W}))\bigg\|_c \\
		&\leq \quad C_\rho\bigg\|g_1(T^Q_{L-1,j}(W_1,\mathds{W})) -
		g_1(T^G_{L-1,j}(W_1,
		\mathds{W}))-\sum_{j\in C^r(x_L)\land j \neq 
			x_L}W^j_R\bigg(g_0(T^Q_{L-1,j}(W_1,W_c)) 
		\bigg)\bigg\|_c \\
		&+d\max_{j\in C^r(x_L)}\bigg\|W^j_R\bigg(g_0(T^Q_{L-1,j}(W_1,\mathds{W})) -
		g_0(T^G_{L-1,j}(W_1,\mathds{W}))\bigg)\bigg\|_c \\
		&\leq \quad C_\rho\bigg\|g_1(T^Q_{L-1,j}(W_1,\mathds{W})) -
		g_1(T^G_{L-1,j}(W_1,
		\mathds{W}))\bigg\|_c+
		\bigg\|-\sum_{j\in C^r(x_L)\land j \neq x_L}W^j_R\bigg(g_0(T^Q_{L-1,j}
		(W_1,W_c)) 
		\bigg)\bigg\|_c \\
		&+d\max_{j\in C^r(x_L)}\bigg\|W^j_R\bigg(g_0(T^Q_{L-1,j}(W_1,\mathds{W})) -
		g_0(T^G_{L-1,j}(W_1,\mathds{W}))\bigg)\bigg\|_c \\
\end{aligned}\end{equation*}

\stab
(a) If none of $Q$ and $G$ contain the center vertices.

\begin{equation*}\begin{aligned}
		\|&R(W_1,\mathds{W},x^Q_L,r)-R(W_1,\mathds{W},x^G_L,r)\|_2 \\
		&=\quad C_\rho\bigg\| - \sum_{j\in C^r(x_L)\land j \neq x_L}
		W^j_R\bigg(g_0(T^G_{L-1,j}(W_1,\mathds{W})) \bigg) \\
		&+\sum_{j\in C^r(x_L)\land j \neq x_L}W^j_R\bigg(g_0(T^Q_{L-1,j}(W_1,
		\mathds{W})) -g_0(T^G_{L-1,j}(W_1,\mathds{W}))\bigg) \\
		&+ \bigg(\mathds{1}(x^Q_L\in C(x^Q_L))g_1(T^Q_{L-1,j}(W_1,\mathds{W})) 
		-
		\mathds{1}(x^G_L\in 
		C(x^G_L))g_1(T^G_{L-1,j}(W_1,\mathds{W}))\bigg)\bigg\|_c \\
		&=\quad C_\rho\bigg\|-\sum_{j\in C^r(x_L)\land j \neq x_L}
		W^j_R\bigg(g_0(T^Q_{L-1,j}(W_1,W_c)) \bigg) \\
		&+\sum_{j\in C^r(x_L)\land j \neq x_L}W^j_R\bigg(g_0(T^Q_{L-1,j}(W_1,
		\mathds{W})) -g_0(T^G_{L-1,j}(W_1,\mathds{W}))\bigg) \bigg\|_c \\
\end{aligned}\end{equation*}


\eat{
%	\clearpage
	
	\begin{equation*}\begin{aligned}
			\|&R(W_1,\mathds{W},x^Q_L,r)-R(W_1,\mathds{W},x^G_L,r)\|_2 \\
			&=\quad C_\rho\bigg\| - \sum_{j\in C^r(x_L)\land j \neq x_L}
			W^j_R\bigg(g_0(T^G_{L-1,j}(W_1,\mathds{W})) \bigg) \\
			&+\sum_{j\in C^r(x_L)\land j \neq x_L}W^j_R\bigg(g_0(T^Q_{L-1,j}(W_1,
			\mathds{W})) -g_0(T^G_{L-1,j}(W_1,\mathds{W}))\bigg) \\
			&+ \bigg(\mathds{1}(x^Q_L\in C(x^Q_L))g_1(T^Q_{L-1,j}(W_1,\mathds{W})) 
			-
			\mathds{1}(x^G_L\in 
			C(x^G_L))g_1(T^G_{L-1,j}(W_1,\mathds{W}))\bigg)\bigg\|_c \\
			&=\quad C_\rho\bigg\|-\sum_{j\in C^r(x_L)\land j \neq x_L}
			W^j_R\bigg(g_0(T^Q_{L-1,j}(W_1,W_c)) \bigg) \\
			&+\sum_{j\in C^r(x_L)\land j \neq x_L}W^j_R\bigg(g_0(T^Q_{L-1,j}(W_1,
			\mathds{W})) -g_0(T^G_{L-1,j}(W_1,\mathds{W}))\bigg) \\
			&+ g_1(T^Q_{L-1,j}(W_1,\mathds{W})) -g_1(T^G_{L-1,j}(W_1,\mathds{W}))
			\bigg\|_c 
			\\
			&=\quad C_\rho\bigg\|g_1(T^Q_{L-1,j}(W_1,\mathds{W})) -g_1(T^G_{L-1,j}
			(W_1,
			\mathds{W}))-\sum_{j\in C^r(x_L)\land j \neq 
				x_L}W^j_R\bigg(g_0(T^Q_{L-1,j}(W_1,W_c)) 
			\bigg) \\
			&+\sum_{j\in C^r(x_L)\land j \neq x_L}W^j_R\bigg(g_0(T^Q_{L-1,j}(W_1,
			\mathds{W})) -g_0(T^G_{L-1,j}(W_1,\mathds{W}))\bigg\|_c \\
			&-\quad C_\rho\bigg\|g_1(T^Q_{L-1,j}(W_1,\mathds{W})) -g_1(T^G_{L-1,j}
			(W_1,
			\mathds{W}))-\sum_{j\in C^r(x_L)\land j \neq 
				x_L}W^j_R\bigg(g_0(T^Q_{L-1,j}(W_1,W_c)) 
			\bigg) \\
			&+\sum_{j\in C^r(x_L)\land j \neq x_L}W^j_R\bigg(g_0(T^Q_{L-1,j}(W_1,
			\mathds{W})) -g_0(T^G_{L-1,j}(W_1,\mathds{W}))\bigg\|_c \\
			&\leq\quad C_\rho\bigg\|g_1(T^Q_{L-1,j}(W_1,\mathds{W})) -g_1(T^G_{L-1,j}(W_1,
			\mathds{W}))-\sum_{j\in C^r(x_L)\land j \neq 
				x_L}W^j_R\bigg(g_0(T^Q_{L-1,j}(W_1,W_c)) 
			\bigg) \\
			&+\sum_{j\in C^r(x_L)\land j \neq x_L}W^j_R\bigg(g_0(T^Q_{L-1,j}(W_1,\mathds{W})) -
			g_0(T^G_{L-1,j}(W_1,\mathds{W}) \\
			&-g_1(T^Q_{L-1,j}(W'_1,W'_2)) -g_1(T^G_{L-1,j}(W'_1,W'_2))-\sum_{j\in C^r(x_L)\land j \neq 
				x_L}W^j_R\bigg(g_0(T^Q_{L-1,j}(W'_1,W'_c)) \bigg) \\
			&+\sum_{j\in C^r(x_L)\land j \neq x_L}W^j_R\bigg(g_0(T^Q_{L-1,j}(W'_1,W'_2)) -
			g_0(T^G_{L-1,j}(W'_1,W'_2))\bigg\|_c \\
	\end{aligned}\end{equation*}
	
	
%	\clearpage
}

\stab
(2) Otherwise, there does not have a match,
that is, $Q$ has ID at the level, but $G$ does not.
\begin{equation*}\begin{aligned}
		\|&R(W_1,\mathds{W},x^Q_L,r)-R(W_1,\mathds{W},x^G_L,r)\|_2 \\
		&= \quad C_\rho\bigg\|-\sum_{j\in C^r(x_L)\land j \neq x_L}
		W^j_R\bigg(g_0(T^G_{L-1,j}(W_1,\mathds{W})) \bigg) \\
		&+\sum_{j\in C^r(x_L)\land j \neq x_L}W^j_R\bigg(g_0(T^Q_{L-1,j}(W_1,
		\mathds{W})) -g_0(T^G_{L-1,j}(W_1,\mathds{W}))\bigg) \\
		&+ \bigg(\mathds{1}(x^Q_L\in C(x^Q_L))g_1(T^Q_{L-1,j}(W_1,\mathds{W})) 
		-
		\mathds{1}(x^G_L\in C(x^G_L))g_1(T^G_{L-1,j}(W_1,\mathds{W}))\bigg)
		\bigg\|_c \\
		&=\quad C_\rho\bigg\|-\sum_{j\in C^r(x_L)\land j \neq x_L}
		W^j_R\bigg(g_0(T^Q_{L-1,j}(W_1,\mathds{W})) \bigg) \\
		&+\sum_{j\in C^r(x_L)\land j \neq x_L}W^j_R\bigg(g_0(T^Q_{L-1,j}(W_1,
		\mathds{W})) -
		g_0(T^G_{L-1,j}(W_1,\mathds{W}))\bigg) \\
		&+g_1(T^Q_{L-1,j}(W_1,\mathds{W}))\bigg\|_c \\
		&=\quad C_\rho\bigg\|g_1(T^Q_{L-1,j}(W_1,\mathds{W}))-\sum_{j\in C^r(x_L)\land j 
			\neq x_L}W^j_R\bigg(g_0(T^Q_{L-1,j}(W_1,\mathds{W})) \bigg) \\
		&+\sum_{j\in C^r(x_L)\land j \neq x_L}W^j_R\bigg(g_0(T^Q_{L-1,j}(W_1,\mathds{W})) -
		g_0(T^G_{L-1,j}(W_1,\mathds{W}))\bigg)\bigg\|_c \\
\end{aligned}\end{equation*}

%\clearpage

\begin{equation*}\begin{aligned}
		&\leq\quad C_\rho\bigg\|\sum_{j\in C^r(x_L)\land j \neq x_L}
		W^j_R\bigg(g_0(T^Q_{L-1,j}(W_1,
		\mathds{W})) \bigg) \\
		&+\sum_{j\in C^r(x_L)\land j \neq x_L}\bigg(g_0(T_{L-1,j}(W_1,\mathds{W})) -
		g_0(T_{L-1,j}
		(W_1',\mathds{W}'))\bigg)\bigg\|_2 \\
		&+ C_\rho\bigg\|\sum_{j\in C^r(x_L)\land j = x_L}\bigg(g_1(T_{L-1,j}(W_1,
		\mathds{W})) -
		g_1(T_{L-1,j}(W_1',\mathds{W}'))\bigg)\bigg\|_2 \\
		&\leq\quad C_\rho C^0_g\sum_{j\in C^r(x_L)\land j \neq x_L}\bigg\|T_{L-1,j}
		(W_1,
		\mathds{W})-T_{L-1,j}(W_1',\mathds{W}')\bigg\| \\
		&+  C_\rho C^1_g\sum_{j\in C^r(x_L)\land j = x_L}
		\bigg\|T_{L-1,j}(W_1,\mathds{W})-T_{L-1,j}(W_1',\mathds{W}')\bigg\|\\
		&=\quad C_\rho \bigg(C^0_g\sum_{j\in C^r(x_L)\land j\neq x_L}\Delta^0_{L-1,j} 
		+ 
		C^1_g\sum_{j\in C^r(x_L)\land j = x_L}\Delta^1_{L-1,j}\bigg)
\end{aligned}\end{equation*}
}



\subsection{Proof of Theorem~\ref{thm-upper-bound}}
Assume by contradiction 
that there  exist a pattern $Q$, a graph $G$,
a vertex $u$ in $Q$ and a vertex $v$ in $G$ 
such that $\HGIN^\kw{v}$ returns \false, 
but \dual returns \true, \ie $(u, v)$ exists
in the relation $S_{(Q, G)}$ computed by \dual.
We deduce a contradiction by analyzing 
the computation of $\HGIN^\kw{v}$.
Observe that
(a) computing
embeddings of $u$ and $v$ can be represented by
trees of height $m$, where $m$ is 
the number of layers in $\HGIN^\kw{v}$
(\cite{Gen-bound}; see Figures~\ref{fig:anti-expressive-0}(a) and~\ref{fig:anti-expressive-0}(b)
for an example);
and (b) vertex $v$ may have more 
adjacent edges in $G$ than $u$, 
\eg paths colored in red in Figure~\ref{fig:anti-expressive-0}(b).
When $\HGIN^\kw{v}$ returns \false, 
the embedding of $v$ is not larger than that of $u$
(see Section~\ref{sec-model}).
Because all functions $\kw{MSG}^{k,r}$ and $\kw{AGG}^k$
$(k\in \mathcal{L})$ in $\HGIN^v$
are injective and monotonic,
there exists a path $p_0$ in the tree representing the computation of 
embedding of $u$ (\eg the path colored in blue
in Figure~\ref{fig:anti-expressive-0}(a)), such that 
no path from $v$ carries the same labels as $p_0$.
However, if this holds, \dual can also identify 
such path and returns \false.
Indeed, the computation of \dual can also be represented
as a tree (see Figures~\ref{fig:anti-expressive-0}(c)
for an example).
The match relation $S_{(Q, G)}$
can be constructed from common parts 
of trees roots at $u$ and $v$ (see Section~\ref{sec-pre}).
When such a path $p_0$ exists in $G$, \dual can also identify 
that $p_0$ is missed in the tree of $v$, 
and $(u, v)$ is not contained in $S_{(Q, G)}$, 
\ie~\dual returns \false, a contradiction.
\looseness = -1

\eat{
Since $\HGIN^v$ also invokes \dual 
to check whether $(u, v)$ exists in $S$, 
$\HGIN^v$ returns \false due to
the prediction function of the model, 
\ie the embedding
of $v$ is not smaller than that of $u$,
from which we can deduce that 
\dual also returns \false, a contradiction.
}





\subsection{Proof of Theorem~\ref{thm-hgin-dual}}
%\begin{proof}
	Consider graphs $G_2$ and $G_3$ in Figure~\ref{fig:difference}, 
	which are extended from~\cite{you2021identity}.
	We can verify that 
	there does not exist a homomorpic mapping 
	from $G_3$ to $G_2$, since $G_3$ contains
	triangles (\ie the subgraph induced by 
	$w_1, w_5$ and $w_6$), but $G_4$ does not have any triangles.
    We next show that
    given $G_3$, $G_2$, $v_1$
    and $w_1$ as input,
     (1) \HGIN returns \false 
    but (2) $(v_1, w_1)\in S_{(G_3,G_2)}$.
    If these hold, then 
    \HGIN is more expressive than \dual,
    \ie \HGIN can distinct more non-homomorphic
    graphs.

\eat{
	\begin{figure}
	%	\vspace{-6.7ex}
	\centerline{\includegraphics[width=0.6\columnwidth]{fig/expressive-1-new.eps}}
	\centering
	\vspace{-1.4ex}
	\caption{Difference of isomorphism and homomorphism} 
	\label{fig:difference-2}
	\vspace{-4ex}
\end{figure}
}
	
	\stab
	(1) We first prove that \HGIN returns \false.
	This is because (a) $v_1$ is contained in a cycle ${\mathcal C}$ of length 3, 
	while $w_1$ does not;
	(b) \HGIN ensures that $\kw{MSG}^{k,r}_1$ is sufficiently 
	larger than $\kw{MSG}^{k,r}_0$; and
	(c) the cycle ${\mathcal C}$ ensures that 
	computing the embedding of $v_1$ 
	invokes more $\kw{MSG}^{k,r}_1$, 
	which leads to that the embedding
	of $v_1$ is larger than that of $w_1$.
	Consider two trees in Figure~\ref{fig:express-4}, 
	which represent the computation of embeddings of 
	$v_1$ and $w_1$. The vertices, which message passing functions
	are $\kw{MSG}^{k,r}_1$, is marked red.
	We can see that
	the computation of embedding of $v_1$ invokes
	more functions $\kw{MSG}^{k,r}_1$.
	\looseness = -1
	
\stab	
	(2) We next show that 
	$(v_1, w_1)\in S_{(G_3,G_2)}$.
	That is, $(v_1, w_1)$ 
	exists in the dual simulation of $G_3$ and $G_2$.
	Observe that both  $G_3$ and $G_2$ are regular, 
	and all vertices have the same number of
	neighborhoods.
	Therefore, all vertices pairs 
	exist in $V_{G_3}\times V_{G_2}$.
%\end{proof}

\eat{
\begin{figure}
	%	\vspace{-6.7ex}
	\centerline{\includegraphics[width=0.6\columnwidth]{fig/expressive-1-1.eps}}
	\centering
	\vspace{-1.4ex}
	\caption{Difference of isomorphism and homomorphism} 
	\label{fig:difference-2}
	\vspace{-4ex}
\end{figure}
}


\begin{figure*}[h]
%	\vspace{-5.7ex}
	\begin{center}
		\begin{minipage}[t]{0.8\textwidth}
			%\hfill
			\subfigure{{\includegraphics[width=0.4\columnwidth]{fig/expressive-1-new.eps}}}
			\hfill
			\subfigure{{\includegraphics[width=0.4\columnwidth]{fig/expressive-4.eps}}}
			\hfill
		\end{minipage}
	\end{center}
	\vspace{-3.4ex}
	\caption{Proof of Theorem~\ref{thm-hgin-dual}} 
    \label{fig:express-4}
	\vspace{-3ex}
\end{figure*}


\eat{
	\begin{figure}[h]
	%	\vspace{-6.7ex}
	\centerline{
		\includegraphics[width=0.4\columnwidth]{fig/expressive-1-1.eps}
		\includegraphics[width=0.3\columnwidth]{fig/expressive-4.eps}}
	\centering
	\vspace{-1.4ex}
	\caption{Proof of Theorem~\ref{thm-hgin-dual}} 
	\label{fig:express-4}
	\vspace{-4ex}
\end{figure}
}

\subsection*{A3. Proof of Theorem~\ref{thm-more-expressive-1}}


\begin{figure}
	%	\vspace{-6.7ex}
	\centerline{\includegraphics[width=0.7\columnwidth]{fig/example.eps}}
	\centering
	\vspace{-1.4ex}
	\caption{Example in Theorem~\ref{thm-more-expressive-1}} 
	\label{fig:label-embedding}
%	\vspace{-4ex}
\end{figure}


\eat{
\begin{theorem}
	When the embeddings are finite domain, 
	then Dualsimulation is more expressive than vanilla
	GNN.
	\label{thm-more-expressive}
\end{theorem}
}

%\proofS
%(2) 
Assume that the embeddings are finite dimension, 
and let the embedding be $d$-dimension.
We select $2d$ distinct edge labels
$l_1, \ldots, l_{2d}$, 
which embeddings are $e_1, \ldots, e_{2d}$, respectively.
%and construct $2d$ vectors to represent these labels.
Since the embeddings are $d$-dimension, 
we can pick $d$ embeddings $e'_1, \ldots, e'_d$
such that the $i$ index of $e'_i$ 
is maximum among all $2d$ embeddings $e_1, \ldots, e_{2d}$.
Let $l'_1, \ldots, l'_d$ be the labels, which embeddings
are $e'_1, \ldots, e'_d$, respectively.
Next we construct a star $Q$ 
with $d$ edges, which labels 
are the picked $d$ labels $l'_1, \ldots, l'_d$.
We can deduce a contradiction by comparing
the embeddings of $Q$ and all the $2d$ labels.
(1) Because $Q$ consists of $d$ labels, 
the embedding of $Q$ must be larger than 
the embeddings of these $d$ labels, based on the 
order-embedding space (see the prediction function
of \HGIN);
(2) each of these $d$ vectors 
has the maximum value for at least one index,
then the embedding of $Q$ must be larger than 
the embeddings of all $2d$ labels,
which is a contradiction, 
since $Q$ has only $d$ labels,
and $l_1, \ldots, l_{2d}$ are distinct.




\eat{
Since the order-embedding space is monotoncity, 

each index of $Q$ is maximum, 
So all $2d$ labels is a subgraphs of $Q$, which is a construction.
}	

Consider the six edge in Figure~\ref{fig:label-embedding}(1),
which carry distinct labels
$l_1, \ldots, l_6$. Their embeddings in two-dimensional space
as shown in Figure~\ref{fig:label-embedding}(b).
Since these labels are distinct, 
none of these labels is homomorphic 
to any of the other labels.
Therefore, given any embeddings
$e_1 = (a_1, a_2)$ and 
$e_2 = (b_1, b_2)$ of two labels 
$l_1$ and $l_2$, respectively,
if $a_1\leq b_1$, then $a_2> b_2$;
and if $a_2\leq b_2$, then $a_1> b_1$ 
(see Figure~\ref{fig:label-embedding}(2)).
Then we pick two labels 
$l_3$ and $l_5$, which have the largest 
values in the second dimension and the first dimension,
respectively,
and construct a pattern $Q_c$ in 
Figure~\ref{fig:label-embedding}(3).
We can verify that 
(I) $Q_c$ is homomorphic
to none of the edges in Figure~\ref{fig:label-embedding}(1).
(II) %On the other hand, 
since $l_3$ and $l_6$ have the largest 
values in the second dimension and the first dimension,
respectively,
the embedding of $Q_c$ is larger than embeddings
of all edges in Figure~\ref{fig:label-embedding}(1),
\ie the six edges in Figure~\ref{fig:label-embedding}(1)
can be homomorphic to $Q_c$,
which contradicts the fact that 
$l_1, \ldots, l_6$ are distinct.


\subsection*{A4. Proof of Theorem~\ref{thm-generalization-bound}}
We establish the Rademacher complexity  of \HFrame
by extending the proof of 
the Rademacher complexity for \GNNs~\cite{Gen-bound}.
Since \dual does not have any parameters, 
and works for any inputs, we only 
bound the Rademacher complexity of \HGIN.
We prove the bound in multiple steps:
(1) represent the computation of \HGIN
by a tree;
(2) quantify the impact of change 
of parameters on the embeddings;
(3) bound the changes in prediction probability;
(4) bound the covering numbers for \HGIN;
and (5) finally bound the Rademacher complexity
using the covering numbers.

In the following, we
show these steps one by one. 
In the proof of Theorem~\ref{thm-upper-bound}, 
the computation of \HGIN is represented by a 
computation tree.
It remains to prove steps (2)-(5).	
	
\stitle{$\bullet$ Step 2: Bound the impact of parameters}.
We first construct a formula 
to represent changes of embeddings, and then bound the changes.


Based on Equations (1)-(4), 
we can deduce the following upper bound
for the changes of embeddings:
\begin{small}
\begin{align*}
\Delta_{L}
  \leq& \quad 2C_{\kw{COM}}\left\|(W_{1}-W_{1}^{\prime})x_{L}\right\|_{2} \\
& +\quad C_{\kw{COM}}\sum_{r\in R}\|(W_rR^+(W_{1},\mathds{W},x_{L})-W_r^{\prime}R^+(W_{1},\mathds{W},x_{L}))\|_{2}\\
& + C_{\kw{COM}}\sum_{r\in R}\|(W^{\prime}_rR^+(W_{1},\mathds{W},x_{L})-W_r^{\prime}R^+(W_{1}^{\prime},\mathds{W}^{\prime},x_{L}))\|_{2}\\
& +\quad C_{\kw{COM}}\sum_{r\in R}\|(W_rR^-(W_{1},\mathds{W},x_{L})-W_r^{\prime}R^-(W_{1},\mathds{W},x_{L}))\|_{2}\\
& + C_{\kw{COM}}\sum_{r\in R}\|(W^{\prime}_rR^-(W_{1},\mathds{W},x_{L})-W_r^{\prime}R^-(W_{1}^{\prime},\mathds{W}^{\prime},x_{L}))\|_{2} \tag{11}
\end{align*}
\end{small}

To bound $\Delta_{L}$, 
we bound the following:
\begin{itemize}
\item[(a)] $\|(W_{1}{-}W_{1}^{\prime})x_{L}\|_{2}$;
\item[(b)]  $\|W_rR^+(W_{1},\mathds{W},x_{L}){-}W_r^{\prime}R^+(W_{1},\mathds{W},x_{L})\|_{2}$;
\item[(c)]  $\|W^{\prime}_rR^+(W_{1},\mathds{W},x_{L}){-}W_r^{\prime}R^+(W_{1}^{\prime},\mathds{W}^{\prime},x_{L})\|_{2}$;
\item[(d)]  $\|W_rR^-(W_{1},\mathds{W},x_{L})-W_r^{\prime}R^-(W_{1},\mathds{W},x_{L})\|_{2}$;
and 
\item[(e)]  $\|(W^{\prime}_rR^-(W_{1},\mathds{W},x_{L}){-}W_r^{\prime}R^-(W_{1}^{\prime},\mathds{W}^{\prime},x_{L})\|_{2}$.
\end{itemize}
%\looseness=-1

For term (a), we can assume an upper bound
for the norm of $W_1$ and $x_L$. 
For (b), 
we have 
$\|W_rR^+(W_{1},\mathds{W},x_{L})-W_r^{\prime}R^+(W_{1},\mathds{W},x_{L})\|_{2}
\leq \|(W_r-W_r^{\prime})R^+(W_{1},\mathds{W},x_{L})\|_{2}$.
Then we only need to bound $\|R^+(W_{1}^{\prime},\mathds{W}^{\prime},x_{L})\|_{2}$.
For (c), 
we have that 
$\|W^{\prime}_rR^+(W_{1},\mathds{W},x_{L})-W_r^{\prime}R^+(W_{1}^{\prime},\mathds{W}^{\prime},x_{L})\|_{2}
\leq \|W^{\prime}_r(R^+(W_{1},\mathds{W},x_{L})-R^+(W_{1}^{\prime},\mathds{W}^{\prime},x_{L}))\|_{2}$.
Therefore, we bound 
$\|R^+(W_{1},\mathds{W},x_{L})-R^+(W_{1}^{\prime},\mathds{W}^{\prime},x_{L})\|_{2}$.
The cases (d) and (e) can be analyzed similarly.
In the following, we only provide
upper bounds for $\|R^+(W_{1}^{\prime},\mathds{W}^{\prime},x_{L})\|_{2}$
and $\|R^+(W_{1},\mathds{W},x_{L})-R^+(W_{1}^{\prime},\mathds{W}^{\prime},x_{L})\|_{2}$.

\vspace{2ex}

\etitle{Upper bound for $\|R^+(W_{1},\mathds{W},x_{L})-R^+(W_{1}^{\prime},\mathds{W}^{\prime},x_{L}))\|_{2}$}.
We start with the upper bounds for term $\|R^+(W_{1},\mathds{W},x_{L})-R^+(W_{1}^{\prime},\mathds{W}^{\prime},x_{L}))\|_{2}$.
Based on the definition of 
$R^+(W_{1},\mathds{W},x_{L})$, we have the following.

\begin{scriptsize}	
%\vspace{-7ex}
\begin{align*}
&\|R^+(W_1,\mathds{W},x_L,r)-R^+(W_1',\mathds{W}',x_L,r)\|_2\\ % \tag{eq-1} 
  &\leq C_\rho C_{\kw{MSG}_0}C_{\kw{MSG}_1}d\max_{j\in N_r^+(x_L)}\Delta_{L-1,j}\tag{12}
\end{align*}
\end{scriptsize}

Equation (12) holds, since we assume that 
the degrees of vertices in $G$ are bounded by 
a constant $d$.
\looseness = -1

Observe that the upper bound for
$\Delta_{L}$ depends on that of $\|R^+(W_1,\mathds{W},x_L,r)-R^+(W_1',\mathds{W}',x_L,r)\|_2$,
which in turn dependents on $\Delta_{L-1,j}$.
Then we deduce a bound for $\|R^+(W_1,\mathds{W},x_L,r)-R^+(W_1',\mathds{W}',x_L,r)\|_2$
by expanding the bounds (see below). 


\etitle{Upper bounds for $\|R^+(W_{1}^{\prime},\mathds{W}^{\prime},x_{L}))\|_{2}$}.
Based on the definition of $\|R^+(W_{1},\mathds{W},x_{L},r)\|_{2}$, 
we can deduce the following.

\begin{align*}
  \|R^+(W_{1},\mathds{W},x_{L},r)\|_{2} 
 & \leq C_\rho C_{\kw{MSG}_0}C_{\kw{MSG}_1}d\max_{j\in N_r^+(x_L)}\bigg\|T_{L-1,j}(W_1,\mathds{W})\bigg\|_2
\end{align*}


The equations hold due to the reasons given above.
Observe that $\|R^+(W_{1},\mathds{W},x_{L},r)\|_{2}$ 
is bounded using $\bigg\|T_{L-1,j}(W_1,\mathds{W})\bigg\|_2$,
\ie the changes of the predictions of neighbors.
We continue to expand $\bigg\|T_{L-1,j}(W_1,\mathds{W})\bigg\|_2$
as follows.

\vfill\null
\begin{align*}
    &\quad\quad\bigg\|T_{L-1,j}(W_1,\mathds{W})\bigg\|_2\\
    & \leq\quad 2C_{\kw{COM}}B_{1}B_{x}+C_{\kw{COM}}|R|B_{2}\bigg\|
    R^+(W_{1},W_{2},x_{L-1,j},r)\bigg\|_{2}\\
    &\quad\quad+C_{\kw{COM}}|R|B_{2}\bigg\|
    R^-(W_{1},W_{2},x_{L-1,j},r)\bigg\|_{2}.  \tag{16}
\end{align*}



And Equation (16) holds, since 
we assume that 
the norms of $W_1$ and $X_i$
is bounded by $B_1$ and $B_x$, respectively.
%
Using the above bound, we bound $\|R^+(W_1,\mathds{W},x_L,r)\|_2$ %can be bounded
as follows.

\begin{align*}
&  \|R^+(W_1,\mathds{W},x_L,r)\|_2 \\
 & \leq \quad 2C_\rho C_{\kw{MSG}_0}C_{\kw{MSG}_1}C_{\kw{COM}} B_1B_xd\frac{(\mathcal{C}d)^L-1}{\mathcal{C}d-1}. \tag{13}
\end{align*}

Let ${\mathcal C} = 2C_\rho C_{\kw{MSG}_0}C_{\kw{MSG}_1}C_{\kw{COM}} B_2 |R| d$.
Denote $M$ as $\frac{\left(\mathcal{C}d\right)^{L}-1}{\mathcal{C}d-1}$.
That is, $\|R^+(W_1,\mathds{W},x_L,r)\|_2  \leq {\mathcal C} B_1  B_xM$.


\etitle{Upper bounds for $\Delta_{L}$}.
Combining Equations (11), (12) and (13), 
we get

\eat{
\begin{gather*}
\Delta_{L}\quad\leq\quad MB_{x}\left\|W_{1}-W_{1}^{\prime}\right\|_{2} 
+ M||R(W_1,\mathds{W},x_L)||_2||\mathds{W}-\mathds{W}^{\prime}||_2.
\end{gather*}
}

\begin{align*}
	\Delta_{L} & \leq\quad 4M \bigg(2B_{x}||W_{1}-W_{1}^{\prime}||_{2}\\
	&\quad+\max_{r\in R}\|W_r-
	W_r^{\prime}\|_2(\|R^+(W_1,\mathds{W},x_L,r)\|_2\\
	&\quad+\|R^-(W_1,\mathds{W},x_L,r)\|_2)\bigg) \tag{14}
\end{align*}


\stitle{$\bullet$ Step 3: Bound the prediction changes}.
We first reform the change of 
the prediction function as follows. 

\begin{equation*}\begin{aligned}
		P_{L} & \triangleq\quad\|T^Q_{L}(W_{1},\mathds{W})-T^G_{L}(W_{1},
		\mathds{W})\|_c \\
		& =\quad\bigg\|\kw{COM}(W_{1}x^Q_{L}+\sum_{r\in R} W_rR^+(W_{1},\mathds{W},x^Q_{L},r)), W_{1}x^Q_{L}\\
		& \quad\quad+\sum_{r\in R} W_rR^-(W_{1},\mathds{W},x^Q_{L},r))\\
			& \quad\quad-\kw{COM}(W_{1}x^G_{L}+\sum_{r\in R} W_rR^+(W_{1},\mathds{W},x^G_{L},r)), W_{1}x^G_{L}\\
			&\quad\quad+\sum_{r\in R} W_rR^-(W_{1},\mathds{W},x^G_{L},r)) \bigg\|_c\\
%		&=W_{1}(x^Q_{L}-x^G_{L}) + \sum_{r\in R} W_r(R(W_{1},
%		\mathds{W},x^Q_{L},r)-R(W_{1},\mathds{W},x^G_{L},r))\bigg\|_c\\
\end{aligned}\end{equation*}

Then the change of prediction can be bounded as follows.

\begin{scriptsize}
\begin{align*}
		\nabla_{L} & \triangleq \quad |P_L - P'_L|\\
&\leq \quad 4M \bigg(2B_{x}||W_{1}-W_{1}^{\prime}||_{2}+M\max_{r\in R}\|W_r-
W_r^{\prime}\|_2(\|R^+(W_1,\mathds{W},x_L,r)\|_2\\
&\quad\quad+\|R^-(W_1,\mathds{W},x_L,r)\|_2)\bigg)\tag{15}
\end{align*}
\end{scriptsize}

Equation (15) holds due to Equation (14) above.

\eat{
Here, we show that $\|a\|_c - \|b\|_c\leq \|a-b\|_c$.
Assume that $\|a\|_c = a_i$ and $\|b\|_c=b_j$.
Then $b_i\leq b_j$ and $-b_i\geq -b_j$. 
Hence, $a_i-b_j\leq a_i-b_i \leq \max_{k\in [1, d]} a_i-b_i$.
That is, $\|a\|_c - \|b\|_c\leq \|a-b\|_c$.
}

\stitle{$\bullet$ Step 4: Bound the covering numbers}.
To bound the change in $\nabla_{L}$ by 
a given threshold $\epsilon$, we use the covering number.
Intuitively, given a set $\M$ of matrices
that have a bound for the norm $\|\cdot\|_F$,
\ie $\|M\|_F\leq \lambda$ for all $M\in \M$,
and a number $\eta$, its covering 
number ${\mathcal N}(\M, \eta, \|\cdot\|_F)$ is defined as the minimum number
of balls of radius $\eta$, needed to cover
all matrices in $\M$~\cite{Covering-number}.
To bound the covering number, 
we use a lemma from~\cite{Covering-1}, 
which states that 
$${\mathcal N}(\M, \eta, \|\cdot\|_F)\leq \bigg(1+\frac{2d\lambda}{\epsilon}\bigg)^{d^2}$$

Here, $\varepsilon$ is a given threshold, $\lambda$
is the bound for the norm of matrices in $\M$, 
and $d$ is the dimension 
of the matrix.

Since the bound in Equation (20) has two kinds
of learned matrices, \ie $W_1$ and matrices in $\mathds{W}$.
Using this lemma, we deduce that 


\begin{align*}\mathcal{N}\left(W_{1},\frac{\epsilon}{8MB_{x}},\|\cdot\|_{F}\right)
\leq\left(1+\frac{16MB_{x}B_{1}\sqrt{d}}{\epsilon}\right)^{d^{2}}\end{align*}

\begin{align*}\mathcal{N}\left(W_r,\frac{\epsilon}{8M{\mathcal R}},
	\|\cdot\|_{F}\right)\leq\left(1+\frac{16M{\mathcal R}B_{2}\sqrt{d}}{\epsilon}
	\right)^{d^{2}} \quad (r\in R)
\end{align*}

Here, ${\mathcal R} = \max(\|R^+(W_1,\mathds{W},x_L,r)\|_2,
\|R^-(W_1,\mathds{W},x_L,r)\|_2) \leq {\mathcal C} B_1  B_xM$.
Therefore, the overall covering number for \HGIN is bounded
by the production of all these covering numbers~\cite{Gen-bound}.
\begin{align*}P=\left(1+\frac{16M\sqrt d\max\{B_xB_1,{\mathcal R}B_2\}}
{\epsilon}\right)^{2|R|d^2}\end{align*}

Because $\varepsilon$ is sufficiently small, 
\eg $\epsilon<4M\sqrt{d}\max\{B_xB_1,{\mathcal R}B_2\}$,
we have that $\log P\leq 2|R|d^2\log\left(\frac{32M\sqrt{d}\max\{B_xB_1,{\mathcal R}
B_2\}}{\epsilon}\right)$.
We will use this upper bound to prove the Rademacher complexity.

\eat{
\begin{equation*}
	3r^2\log\left(1+\frac{6M\sqrt{r}\max\{B_xB_1,\overline{R}B_2\}}{\epsilon}\right)
\end{equation*}
\begin{equation*}
	\epsilon<4 M\sqrt{r}\max\{B_xB_1,\overline{R}B_2\}
\end{equation*}
\begin{equation*}
	3r^2\log\left(\frac{8M\sqrt r\max\{B_xB_1,\overline{R}B_2\}}{\epsilon}\right)
\end{equation*}
}

\stitle{$\bullet$ Step 5: Bound for the Rademacher complexity}.
We finally can bound the Rademacher complexity.
To this end, we also apply a lemma from~\cite{RB-1},
which states that 
given a function ${\mathcal F}$ taking values
in $[0,1]$, 
its Rademacher complexity 
$\hat{\mathcal{R}}$ 
is bounded by $\inf_{\alpha>0}\left(\frac{4\alpha}{\sqrt{N_T}}+\frac{12}{N_T}\int_{\alpha}^{\sqrt{N_T}}\sqrt{\log\mathcal{N}(\mathcal{F},\epsilon,\|\cdot\|_F)}d\epsilon\right)$,
where $\mathcal{N}(\mathcal{F},\epsilon,\|\cdot\|_F)$ is the covering 
number of $\mathcal{F}$, and
$N_T$ is the number of training data.


From Step-4, we know that the covering number ${\mathcal N}$
is bounded by $\left(1+\frac{16M\sqrt d\max\{B_xB_1,{\mathcal R}B_2\}}
{\epsilon}\right)^{2|R|d^2}$.

Therefore, we can bound the covering number as follows.

\begin{align*}
\log \mathcal{N}(\mathcal{F},\epsilon,\|\cdot\|_F)&\leq 
\log \mathcal{N}(P,\frac{\epsilon}{2},\|\cdot\|_F)\\
&
\leq 2|R|d^2\log\left(\frac{32M\sqrt{d}\max\{B_xB_1,\overline{R}
	B_2\}}{\epsilon}\right)
\end{align*}

Setting $\alpha  = \frac{2}{\sqrt{N_T}}$, we have that 
$\hat{\mathcal{R}}(\mathcal{T})\leq\frac{8}{N_T}+
\frac{12}{\sqrt{N_T}}\sqrt{2|R|d^2\log T}$,
where $T=16M\sqrt{{N_T}d}\max\{B_xB_1,{\mathcal R}B_2\}$,
and $\mathcal{I}$ is the class of functions
mapping the computation of embedding to the prediction function $\nabla_{L}$.

We can verify that the max margin loss function 
is 1-Lipschitz.
Then ${\hat{\mathcal{R}}(\mathcal{J}_\gamma)}
%\leq \hat{\mathcal{R}}
%_T(\mathcal{I})
\leq\frac{8}{N_T}+\frac{12d}{\sqrt{N_T}}
\sqrt{2|R|d^2|R|\log T}$.

\stitle{The 1-Lipschitz of the loss function}.
We next show that the max margin loss function 
is 1-Lipschitz, 
\ie $|{\mathcal L}(h_u,h_v)-{\mathcal L}(h'_u,h'_v)|\leq 
|{\mathcal M}(h_u,h_v)-{\mathcal M}(h'_u,h'_v)|$.
Since the loss functions for positive data and negative data
are different,
we consider all possible four cases:
(1) both $(h_u,h_v)$ and $(h'_u,h'_v)$ are positive;
(2) both $(h_u,h_v)$ and $(h'_u,h'_v)$ are negative;
(3) $(h_u,h_v)$ is positive and $(h'_u,h'_v)$ is negative; and
(4) $(h_u,h_v)$ is negative and $(h'_u,h'_v)$ are positive.



\stab
(1) If both $(h_u,h_v)$ and $(h'_u,h'_v)$ are positive, 
then their loss are $-\alpha$.
Hence, 
$|{\mathcal L}(h_u,h_v)-{\mathcal L}(h'_u,h'_v)| = 0\leq 
|{\mathcal M}(h_u,h_v)-{\mathcal M}(h'_u,h'_v)|$.

\eat{
\stab
(2) If both $(Q_1,G_1)$ and $(Q_2,G_2)$ are negative,
and the margins between $(Q_1,G_1)$ and $(Q_2,G_2)$ are smaller than
$\alpha$, then the inequality holds.

\stab
(3) If both $(Q_1,G_1)$ and $(Q_2,G_2)$ are negative,
and the margins between $(Q_1,G_1)$ and $(Q_2,G_2)$ are larger than 
$\alpha$, then the loss is 0, and the inequality holds.
}

\stab
(2) If both $(h_u,h_v)$ and $(h'_u,h'_v)$ are negative,
when both ${\mathcal M}(h_u,h_v)$
and ${\mathcal M}(h'_u,h'_v)$ are not larger than 
$\alpha$, we have that 

\begin{align*}
&{\mathcal L}(h_u,h_v)-{\mathcal L}(h'_u,h'_v)|\\
&=(\alpha -\|\max(0,h_u-h_v)\|^2_2) - (\alpha - \|\max(0,h'_u-h'_v)\|^2_2)\\
&=-\|h_u-h_v\|^2_2 + \|h'_u-h'_v\|^2_2 \\
&\leq |{\mathcal M}(h_u,h_v)-{\mathcal M}(h'_u,h'_v)| \tag{16}
\end{align*}

Equation (21) holds, since 
$\mathcal{M}(h_u, h_v) = \cdd{\max(0, h_u-h_v)}_2^2$.

\vspace{2ex}
When ${\mathcal M}(h_u,h_v)\geq \alpha$ or
 ${\mathcal M}(h'_u,h'_v)\geq  \alpha$,
 the proof is similar.


\stab
(3) If $(h_u,h_v)$ is positive and $(h'_u,h'_v)$ is negative,
and ${\mathcal M}(h'_u,h'_v)\leq \alpha$, we have the following

$\|\max(0,h_u - h_v)\|_2 - \alpha  - (\alpha - \|\max(0,h'_u - h'_v)\|)\\
=   \|h'_u - h'_v \|_2- 2\alpha\\
\leq \|h_u - h_v\|_2 +  \|h'_u - h'_v\|_2- 2\|h'_u - h'_v \|_2\\
\leq \|h_u - h_v\|_2 - \|h'_u - h'_v \|_2\\
\leq |{\mathcal M}(h_u,h_v)-{\mathcal M}(h'_u,h'_v)|$





\stab
(4) If $(h_u,h_v)$ is positive and $(h'_u,h'_v)$ is negative,
and ${\mathcal M}(h'_u,h'_v)\geq \alpha$, then 
we have 

$
\|\max(0,h_u - h_v)\|_2 - \alpha  - \max(0,\alpha - \|\max(0,h'_u - h'_v)\|)\\
= - \alpha\\
\leq |{\mathcal M}(h_u,h_v)-{\mathcal M}(h'_u,h'_v)|$


% \clearpage
\subsection*{A5. Complex queries}

\begin{figure}[tb!]
	%	\vspace{-6.7ex}
	\centerline{\includegraphics[width=0.8\columnwidth]{fig/query-example-2.eps}}
	\centering
	\vspace{-1.4ex}
	\caption{Example of complex query}
	\label{fig:query}
	\vspace{-1.4ex}
\end{figure}

When patterns are complex (\eg with multiple cycles
and large number of vertices),
the exact algorithms can be much slow 
and \HFrame is significantly faster than 
them. 
Consider a complex query extracted from Citeseerx,
which is depicted in 
Figure~\ref{fig:query}.
\kw{SIM}-\kw{TD}, \kw{EJ} and \kw{PJ} take
2.691s, 11.164s and 13.472s respectively to conduct the query,
while \HFrame takes only 0.081s.

